\section{Boundary rank and Higgs tails}\label{sec:boundary-rank}

This section proves the boundary rank loss and the weighted Higgs-tail
comparison that turn the marked-period construction into adjoint
bigness and, later, base nonvanishing.

We retain the marked data of \Cref{sec:marked-periods}.  In particular,
$G$ is the connected rational adjoint monodromy group, $\mathbb A$ is an
irreducible representation of $G_{\C}$ with the compatible polarized
Hodge grading, and $E_0$ is its whole highest Hodge step.  The marks are
$j=(E_0,j_{\rm u})$.  Their orbit $J$ is a product of projective flag
varieties, and each rational simple factor of $G$ has at least one
complex simple factor acting nontrivially on $J$.  The compact marks
occur on factors with an invariant positive Hermitian metric.  On an
appropriate smooth connected unipotent level the line $\mathcal H$ is
$\det E_0+\mathcal H_{\rm u}$; it is understood downstairs as a rational
line.  Its curvature rank is the generic rank of $dj$, whereas the
quotient $S$ remembers the full Lie period and the compact marks.
These two ranks need not be identified in advance.

\subsection{The flat-connection estimate}

To compare numerical ranks, let $T$ be a local fiber of a flag map in
flat coordinates and let $Z$ be its algebraic closure.  Frames in which
the flag has its fixed value form an algebraic incidence over $Z$.
Its fiber dimension is the stabilizer dimension, whereas it contains
the horizontal lift of $T$.  The following estimate bounds the closure
of that lift from below.  In the orbit calculation these two dimensions
will force the differential on $Z$ to fill the monodromy orbit.

The estimate is the semisimple case of the connection Ax--Schanuel
framework of \cite[Theorem~A, Theorem~3.6, and Remark~3.7]{BCFN}.
We give the argument, following the connection-form proof in
\cite[\S3.1]{BCFN}.

\begin{lemma}\label{lem:connection-dimension}
Let $p:\mathcal P\to Z$ be an algebraic principal bundle for a connected
semisimple complex algebraic group $H$, where $Z$ is smooth and
irreducible.  Suppose $\mathcal P$ carries an algebraic flat principal
connection whose monodromy is Zariski dense in $H$.  Let $U$ be an
irreducible analytic germ contained in a horizontal leaf, and suppose
$p(U)$ is Zariski dense in $Z$.  If $R$ is the Zariski closure of $U$ in
$\mathcal P$, then
\begin{equation}\label{eq:prelim-horizontal-dimension}
 \dim R\geq\dim U+\dim H.
\end{equation}
\end{lemma}

\begin{proof}
Put $\mathfrak h=\Lie H$, and let
$\xi\in H^0(\mathcal P,\Omega^1_{\mathcal P}\otimes\mathfrak h)$
be the absolute connection form.  It has kernel the horizontal
distribution, is the identity on fundamental vertical vectors, and
satisfies
\begin{equation}\label{eq:prelim-maurer-cartan}
 d\xi+\tfrac12[\xi,\xi]=0,\qquad
 R_g^*\xi=\operatorname{Ad}(g^{-1})\xi\quad(g\in H).
\end{equation}
Restrict to a nonempty smooth open subset $R^\circ\subset R$ on which
$\xi|_{TR}$ has constant rank $r$.  The germ $U$ meets this subset, because
it is Zariski dense in $R$.  We may replace $U$ by a small connected
smooth germ there without changing its Zariski closure.  Set
$\mathcal K=\ker(\xi|_{TR^\circ})$, and consider the algebraic
Grassmann map
\[
 \beta:R^\circ\longrightarrow\operatorname{Gr}(r,\mathfrak h),
 \qquad x\longmapsto\xi(T_xR).
\]
For a local holomorphic vector field $V$ in $\mathcal K$, Cartan's
formula and \eqref{eq:prelim-maurer-cartan} give
\[
 \mathcal L_V(\xi|_{R^\circ})
 =\iota_Vd\xi+d(\iota_V\xi)=0.
\]
The local flow of $V$ consequently preserves both the restricted
connection form and its image subspace; hence $d\beta(V)=0$.
Since $TU\subset\mathcal K|_U$, the map $\beta$ is constant on $U$.
Its algebraicity and the Zariski density of $U$ make it constant on
$R^\circ$.  Denote its value by $\mathfrak b\subseteq\mathfrak h$.

This fixed subspace is a Lie subalgebra.  In fact, for $b_1,b_2\in
\mathfrak b$, locally choose tangent fields $V_1,V_2$ on $R^\circ$ with
$\xi(V_i)=b_i$.  Such lifts exist because $\xi|_{TR^\circ}$ is a
surjection onto the constant bundle $\mathfrak b$.  Evaluating
\eqref{eq:prelim-maurer-cartan} yields
\[
 \xi([V_1,V_2])=[b_1,b_2],
\]
so $[b_1,b_2]\in\mathfrak b$.

Suppose $\mathfrak b\ne\mathfrak h$.  Let $A$ be the connected analytic
subgroup of $H$ with Lie algebra $\mathfrak b$, and let $B$ be its
Zariski closure.  Then $B$ is a connected algebraic subgroup and is
proper in $H$.  To prove the latter point, if $B=H$, the fact that $A$
preserves $\mathfrak b$ under its adjoint action would imply that $H$
preserves $\mathfrak b$: the stabilizer of this subspace is algebraic.
Thus $\mathfrak b$ would be an ideal of the semisimple Lie algebra
$\mathfrak h$.  Every such ideal is a sum of simple factors and
integrates to a closed connected normal algebraic subgroup of $H$.
That subgroup is $A$, by uniqueness of a connected analytic subgroup
with a given Lie algebra.  A proper $\mathfrak b$ therefore cannot have
Zariski-dense $A$, a contradiction.

Consider the irreducible algebraic subset
\[
 R_B=\overline{R\cdot B}\subseteq\mathcal P.
\]
It dominates $Z$, since $p(U)$ is dense, and it is invariant under the
right action of $B$.  The equivariance in
\eqref{eq:prelim-maurer-cartan} shows that, at a general smooth point,
\begin{equation}\label{eq:prelim-saturated-tangents}
 \xi(T R_B)\subseteq\Lie B.
\end{equation}
Indeed, on the image of the multiplication map $R^\circ\times B\to
\mathcal P$, tangent vectors coming from $R^\circ$ have connection
values in $\operatorname{Ad}(B)\mathfrak b\subseteq\Lie B$, and
tangent vectors coming from $B$ have values in $\Lie B$.  Generic
smoothness of this dominant map onto $R_B$ proves
\eqref{eq:prelim-saturated-tangents}.

The $B$-orbits in a fiber have dimension $\dim B$.  Domination of $Z$
therefore gives
$\dim R_B\geq\dim Z+\dim B$.  Conversely,
\eqref{eq:prelim-saturated-tangents} and
$\dim\ker\xi=\dim Z$ give the reverse inequality.  Equality holds,
and at general smooth points
\[
 TR_B=\xi^{-1}(\Lie B).
\]
In particular the whole horizontal distribution is tangent to $R_B$.
This is an algebraic tangency statement on all of $R_B$: if a local
algebraic horizontal vector field differentiates an equation of $R_B$,
the result vanishes on its dense smooth locus and hence on $R_B$.
Thus the ideal sheaf of $R_B$ is preserved by horizontal vector fields.
Their local holomorphic flows preserve $R_B$ in both directions.

Choose a nonempty smooth open $Z_0\subseteq Z$ over which the general
fiber of $R_B\to Z$ is nonempty of dimension $\dim B$.  Horizontal
transport along a path in $Z_0$ is defined throughout the path: locally
it is the fundamental solution of a holomorphic linear system after a
faithful representation of $H$.  The preceding ideal invariance shows
that it preserves $R_B$.  In particular the monodromy of $Z_0$ preserves
the nonempty closed subset $(R_B)_z$ of the $H$-torsor $\mathcal P_z$.
Deleting a proper algebraic subset of a smooth complex variety does not
decrease the original monodromy image: the map
$\pi_1(Z_0,z)\to\pi_1(Z,z)$ is surjective, since loops can be perturbed
off a subset of real codimension at least two.  This monodromy is
therefore Zariski dense in $H$.  A closed subset of an $H$-torsor
invariant under a Zariski-dense translation subgroup is invariant under
$H$, and is either empty or the whole torsor.  But $(R_B)_z$ is nonempty
and has dimension $\dim B<\dim H$.  This contradiction proves
$\mathfrak b=\mathfrak h$.

Finally the generic kernel of $\xi|_{TR}$ contains $TU$ and hence has
dimension at least $\dim U$.  Its generic image has dimension $\dim H$,
so rank--nullity gives \eqref{eq:prelim-horizontal-dimension}.
\end{proof}


\subsection{The orbit calculation}

In this subsection, a differential of a period, flag, or projective
line is computed in local flat frames.  A local fiber is a connected
smooth analytic fiber germ on the locus of constant maximal rank.
The phrase Hodge-generic refers to the integral rational Lie variation,
including its tensor constructions, as in \Cref{lem:marked-monodromy}.

\begin{lemma}[Orbit rank]\label{lem:orbit-rank}
Let $B^\circ$ be a smooth connected algebraic base carrying the marked
data above, with Zariski-dense connected monodromy $G$.  The same
conclusions hold for their pullback by a dominant algebraic map with
connected generic fiber.  Consider either the joint flag map $j$, or
an algebraic projective line section $\ell$ of a flat representation
of $G$.

\emph{Local orbit filling.}
Let $T$ be a local fiber at a Hodge-generic point and let $Z$ be its
irreducible Zariski closure.  After passing to a finite level on a
smooth open of $Z$, its connected monodromy is a rational normal
factor $H$ of $G$.  Write $G=H\times Q$.  The full $Q$ period is
constant on the lifted $Z$.  In the flag case the $Q$ compact marks
are constant there as well.  At suitable smooth generic points of
$T$, writing $a$ for $j$ or $\ell$, one has
\begin{equation}\label{eq:boundary-orbit-fill}
 \dim Z-\dim T=\dim(H_{\C}\cdot a),\qquad
 da(T_zZ)=T_{a(z)}(H_{\C}\cdot a(z)).
\end{equation}
In the flag case write $J=J_H\times J_Q$ for the products of its $H$
and $Q$ flag factors, and $j=(j_H,j_Q)$ for the corresponding
projections in local flat frames.  The orbit in
\eqref{eq:boundary-orbit-fill} is then $J_H$.

\emph{Flag ranks and lines.}
For the flag map on a level of $S$, let $q:B\to S_Q$ denote a resolved
connected-fiber quotient for the full $Q$ period together with its
compact marks, and set $s_Q=\dim S_Q$ and $d_H=\dim J_H$.  If
$r=\rk(dj)$, then
\begin{equation}\label{eq:boundary-rank-split}
 \rk(dj_Q)=s_Q,\qquad r=d_H+s_Q.
\end{equation}
On suitable smooth projective models, with $\rho:B\to S$ the level
map, the nef lines split as
\begin{equation}\label{eq:boundary-line-split}
 \rho^*\mathcal H=\mathcal H_H+\mathcal H_Q,
 \qquad \nu(\mathcal H_H)=d_H,
 \qquad \mathcal H_Q=q^*\underline{\mathcal H}_Q,
 \qquad \nu(\underline{\mathcal H}_Q)=s_Q.
\end{equation}
Here all line equalities are rational line-bundle equalities.
\end{lemma}

\begin{proof}
We may choose the initial point outside the countably many exceptional
tensor loci and in the maximum-rank loci of the maps associated with
every rational normal-factor subset of $G$.  There are only finitely
many such subsets.  The restriction to $Z$ has the same generic
Mumford--Tate group as the ambient variation: a tensor condition
holding on $Z$ holds at the chosen Hodge-generic point.  The connected
monodromy normality theorem, in the pure polarized integral setting
specified in \Cref{lem:marked-monodromy}, makes $H$ normal in that
Mumford--Tate group.  Since $H\subseteq G$ and $G$ is semisimple
adjoint, $H$ is a product of rational simple factors and has the
complementary factor $Q$.

The restricted $Q$ monodromy is finite.  Kill it by a finite cover of
a smooth open of $Z$.  The fixed-part theorem then says that the full
$Q$ variation is constant.  Its flat frames are algebraic frames:
the algebraic connection has regular singularities and its trivial
local system is the trivial regular-singular connection.  The
$Q$ compact marks are algebraic in these frames.  They are constant
on the lifted $T$, since $j$ was constant there, and hence on the
lifted $Z$.  Indeed, a lift of a nonempty open subgerm of $T$ is
Zariski dense in an irreducible component of this finite cover.  Its
closure has dimension at least $\dim Z$, because its finite image
contains a Zariski-dense subset of $Z$; an irreducible component of
the cover also has dimension $\dim Z$.  This proves the required
density assertion after finite levels.  In the line case we only
need the constancy of the $Q$ period.

Use the algebraic flat $H_{\C}$-frame torsor over this smooth open of
$Z$.  It is the reduction specified by the parallel tensor
invariants; its connection has Zariski-dense monodromy $H$.  Set
$a_0=a(T)$ in a flat trivialization along $T$.  Consider the
algebraic incidence consisting of frames in which $a$ has value
$a_0$.  Over a point where it is nonempty its fiber is a coset of
the stabilizer of $a_0$ in $H_{\C}$.  Its image contains $T$ and is
constructible, so contains a dense open of $Z$.  If
$d=\dim(H_{\C}\cdot a_0)$, the incidence over that open therefore
has dimension
\[
                  \dim Z+\dim H-d.
\]
It contains the horizontal lift of $T$, whose Zariski closure therefore
has at most this dimension.  Applying
\Cref{lem:connection-dimension} to that horizontal germ gives
\[
 \dim T+\dim H\leq\dim Z+\dim H-d,
 \quad\hbox{or equivalently}\quad
 \dim Z-\dim T\geq d.
\]
This use of the incidence works equally for a locally closed
projective-line orbit; projectivity of the orbit is not needed.

Move within $T$ to a smooth point of $Z$ lying over the dense open
just obtained.  Nonempty open subgerms are still Zariski dense in
$Z$.  In flat frames the values of $a$ on a neighborhood in $Z$
lie in $H_{\C}\cdot a_0$.  Moreover
\[
 \ker(da|_{T_zZ})=T_zZ\cap\ker da=T_zT,
\]
because the ambient map has constant rank near $T$.  Thus
$\dim Z-\dim T\leq d$.  Equality proves
\eqref{eq:boundary-orbit-fill}.  For the joint marks, the
factor-by-factor description of $J$ identifies this orbit with
$J_H$.

We now work with $a=j$ on a level of $S$.  The $Q$ joint data is
constant on $Z$, so $T_zZ\subseteq\ker dq$, while
$\ker dj=T_zT\subseteq\ker dq$ at the chosen points.  The orbit
calculation gives
\[
                   dj_H(\ker dq)=T_{j_H(z)}J_H.
\]
Since $j_Q$ descends through $q$, one has
$\ker dq\subseteq\ker dj_Q$.  Conversely, if $dj_Q(v)=0$, choose
$w\in\ker dq$ with $dj_H(w)=dj_H(v)$.  Then $dj(v-w)=0$, hence
$v-w\in\ker dq$, and therefore $v\in\ker dq$.  It follows that
$\ker dj_Q=\ker dq$.  Computing the rank first along $\ker dq$
and then on its quotient proves \eqref{eq:boundary-rank-split}.
Our initial choice in the maximum-rank loci makes these generic
rank identities, even though the argument was carried out at
points of $T$.

The representation decomposes at connected level as
$\mathbb A=A_H\otimes A_Q$, and its highest step as
$E_0=E_H\otimes E_Q$.  The factors are the compatible Hodge
representations constructed from the separate adjoint Lie tensor
systems; a constant scalar shift of a grading has no effect on
the following determinant calculation.  If $e_H=\rk E_H$ and
$e_Q=\rk E_Q$, then
\[
 \det E_0=(\det E_H)^{\otimes e_Q}
                    \otimes(\det E_Q)^{\otimes e_H}.
\]
Add the compact-mark line on the appropriate factors to define
$\mathcal H_H$ and $\mathcal H_Q$.  Their semipositive curvature
forms have kernels $\ker dj_H$ and $\ker dj_Q$.  The first has
rank $d_H$, by orbit filling and the upper bound $\dim J_H$;
the second has rank $s_Q$.  \Cref{lem:log-metric} gives their
numerical dimensions.

The $Q$ representation and flags descend through $q$ by
\Cref{lem:period-descent}, applied to precisely the rational
normal factors in $Q$.  Its integral lattice inherits the
congruence-level condition, so its quasi-unipotent boundary
monodromies are already unipotent.  Choose smooth projective
models resolving $q$ and all the flag sections.  Unipotent
canonical extension commutes with this pullback: in a monomial
boundary chart the new logarithmic residues are integral sums
of the commuting old nilpotent residues, and the untwisted
filtration pulls back.  Thus the equality of the $Q$ lines on
the open set extends as the equality in
\eqref{eq:boundary-line-split}.  The metric rank on the quotient
is $s_Q$, so the same metric lemma gives
$\nu(\underline{\mathcal H}_Q)=s_Q$.
\end{proof}

\subsection{Strict loss of numerical dimension at the boundary}

The reduced divisor $D_0$ on $S$ consists of the boundary components
whose monodromy on the retained rational Lie variation
$\mathfrak g$ is not the identity.  In particular finite
nonidentity monodromy is included.

\begin{proposition}[Boundary drop]\label{prop:boundary-drop}
For every component $D_i$ of $D_0$,
\begin{equation}\label{eq:boundary-strict-drop}
                 \nu(\mathcal H|_{D_i})<\nu(\mathcal H).
\end{equation}
If $\nu(\mathcal H)=0$, then $D_0$ is empty.
\end{proposition}

\begin{proof}
Put $r=\nu(\mathcal H)=\rk dj$ and use
\Cref{lem:orbit-rank}.  Write $B$ for its smooth unipotent level,
and put $n=\dim B=\dim S$.  On $B$ choose a
smooth boundary component $D'$ mapping generically finitely onto
$D_i$.  Further resolutions and restrictions to the strict
transform are allowed.  For any nef line $L$ on a smooth
projective variety, restriction to a divisor cannot increase
its numerical dimension: choose an ample $A$ with $mA-D'$
effective and compare the nonnegative intersections
$L^kD'A^{n-k-1}$ and $mL^kA^{n-k}$.
Consequently
\[
                 \nu(\mathcal H_H|_{D'})\leq d_H.
\]
If $D'$ does not dominate $S_Q$, then
\[
 \nu(\mathcal H_Q|_{D'})\leq\dim q(D')<s_Q.
\]
The binomial expansion of intersections of nef classes gives
$\nu(P+Q)\leq\nu(P)+\nu(Q)$ for nef $P,Q$.  Applied on $D'$,
this proves the strict inequality in this case.  We may therefore
assume that $D'$ dominates $S_Q$.  At its generic point the
$Q$ system is pulled back from the interior of $S_Q$.

\smallskip\noindent
\emph{Infinite local monodromy.}
The logarithm $N$ of a sufficiently divisible power of the
original monodromy is then nonzero and belongs to the rational
Lie algebra of $H$.  Its $Q$ component is zero because the
$Q$ variation extends from the interior at the generic point
of $D'$.  It acts nontrivially on the $H$ orbit of $E_H$.
To check this last assertion, a nonzero rational element in
a rational simple factor has nonzero projection to every
Galois-conjugate complex simple factor.  A unipotent isometry
of a positive Hermitian space is the identity.  Thus a rational
factor detected only by compact marks cannot support a
nonzero component of $N$.  Every other retained rational factor
has a visible first-kind flag factor, and the action of a
complex simple adjoint group on a nontrivial flag variety is
faithful.  A nonzero component of $N$ therefore moves such a
flag.

Let $p_H$ be the highest Hodge index of $A_H$ and set $p_0=e_Hp_H$.
There is no restriction on the rank $e_H$ of $E_H$.  Form
\[
           \mathbb B=\bigwedge^{e_H}A_H,
           \qquad L_H=\det E_H=F^{p_0}\mathbb B,
           \qquad F^{p_0+1}\mathbb B=0.
\]
The highest filtration step of $\mathbb B$ is exactly the
one-dimensional line $L_H$.  Denote the weight of $\mathbb B$
by $w$.  On the open stratum of $D'$, untwist by
$\exp(-\log(z)N/(2\pi\sqrt{-1}))$ in a local normal coordinate
$z$ and let $[v]$ be the limiting highest line.  Schmid's
nilpotent-orbit theorem and limiting mixed Hodge theorem
\cite[Theorems~4.12 and~6.16]{Schmid} apply to the polarized
unipotent variation.  The same assertions hold in $\mathbb B$:
it is a Hodge-compatible summand of rational adjoint tensor
systems, whose flat projectors commute with $N$ and preserve
the two limiting Hodge filtrations and the weight filtration.

Write $W_\bullet=W(N)_\bullet$, centered at $w$.  The line
$F^{p_0}$ occupies a single summand $I^{p_0,q}$ of the Deligne
splitting, since there are no summands of first index greater
than $p_0$.  Set $p_0+q=w+k$.  The primitive decomposition of
the monodromy weight grades shows that $k\geq0$ and that the
class of $v$ in $\Gr^W_{w+k}$ is primitive.  Here is the
specific reason.  A nonprimitive term of first Hodge index
$p_0$ has the form $N^j u$ with $j>0$ and with $u$ of first
index $p_0+j$, which is impossible.  Likewise a weight grade
below $w$ is entirely obtained from positive powers of $N$
on higher primitive grades, and so cannot contain the highest
line.  Monodromy Lefschetz gives a nonzero $N^k[v]$ in
$\Gr^W_{w-k}$ and zero $N^{k+1}[v]$.  In fact
\begin{equation}\label{eq:boundary-actual-primitive}
                         N^kv\ne0,\qquad N^{k+1}v=0
\end{equation}
for the actual vector: $N$ has bidegree $(-1,-1)$ in the
Deligne splitting, so each $N^jv$ lies in a single summand;
its vanishing on the corresponding weight grade is its
vanishing as a vector.  The primitive decomposition and
Lefschetz isomorphisms used here are those of the monodromy
weight filtration, as in \cite[Lemma~6.4 and Theorem~6.16]{Schmid}.

Let $O_H$ be the closed orbit of the highest line of
$\mathbb B$, namely the exterior-power embedding of the
$H$ orbit of $E_H$.  The untwisting acts through $H_{\C}$,
so $[v]\in O_H$.  By \eqref{eq:boundary-actual-primitive},
\[
             \lim_{\tau\to\infty}\exp(\tau N)[v]=[N^kv]
                         \in O_H\cap\mathbf P(\ker N).
\]
The intersection on the right is a proper closed subset of
the irreducible variety $O_H$, because $N$ acts nontrivially
on $O_H$.  Project $N^kv\in W_{w-k}$ to $\Gr^W_{w-k}$ and
then apply the fixed inverse of
\[
                N^k:\Gr^W_{w+k}\xrightarrow{\ \sim\ }
                                      \Gr^W_{w-k}.
\]
The resulting projective point is the highest graded line
$[v]_{\Gr}$.  These maps are fixed linear maps in local flat
limiting coordinates.  Projection is taken on the open set
where its value is nonzero; this set contains all the points
under consideration by Lefschetz.  The local variation rank
of the highest graded line is therefore at most
$\dim O_H-1$.  Together with all the $H$ compact marks its
rank is at most
\begin{equation}\label{eq:boundary-graded-rank}
                    \dim O_H-1+\dim J_{{\rm u},H}=d_H-1.
\end{equation}

\smallskip\noindent
\emph{From limiting rank to numerical dimension.}
\phantomsection\label{proof:limiting-weight-comparison}
We next justify that the graded rank just obtained bounds the numerical
dimension of the restricted line, including at crossings.
One cannot simply restrict the original singular Hodge
metric to $D'$.  Instead, the pure weight grades of the
limiting mixed Hodge structures give polarized pure
variations on the open stratum of $D'$
\cite[Proposition~2.10]{CK}.  Their filtrations
are holomorphic and transverse, by the extended filtration
and the tangential logarithmic connection.  Their primitive
parts have the monodromy polarizations, and their remaining
parts are polarized through the Lefschetz decomposition.

There is a canonical extension description on the whole
component.  In a Deligne frame at a crossing
$z_1\cdots z_a=0$, with $D'=(z_1=0)$, the transverse residue
is the constant nilpotent $N_1$ and the remaining residues
$N_2,\ldots,N_a$ commute with it.  The filtration $W(N_1)$
therefore consists of subbundles of the restricted extended
flat bundle.  Every remaining residue preserves this
filtration and induces a nilpotent operator on each weight
grade.  Changing the normal coordinate or the tangential
lift changes the tangential connection by a multiple of
$N_1$; this acts trivially on the weight grades.  Their flat connections are thus well defined. The remaining logarithmic
exponentials are exactly the Deligne extensions of these flat graded
variations. Separately, each polarized pure graded variation has its
Schmid-extended Hodge subbundles inside that canonical flat bundle.
We do not identify a possibly jumping intersection with a weight step
with an extended Hodge subbundle. Instead we use the independent Hodge
extension, and below map into it by continuation from the open stratum.
This distinction proves the needed comparison also at crossings.

On the stratum, $L_H|_{D'}$ lies in $W_{w+k}$ and maps
nontrivially to its highest graded line.  Inclusion in the
subbundle $W_{w+k}$ extends everywhere because it holds
generically.  Projection to $\Gr^W_{w+k}$ is regular.  The highest line of the graded variation extends as its independent
Schmid Hodge subbundle. The projection lands in it: its image in the
locally free quotient of the canonical flat graded bundle vanishes on
the dense open stratum and hence vanishes everywhere. The resulting
map can have zeros at further boundary strata; these produce an
effective difference in the direction needed below.  If necessary resolve
the additional algebraic filtration-rank loci on $D'$.
The loops newly omitted from the interior have trivial
monodromy, and all boundary monodromies remain unipotent;
the preceding description is preserved on this resolution.
We obtain a generically nonzero morphism of line bundles
\begin{equation}\label{eq:boundary-line-comparison}
                 L_H|_{D'}\longrightarrow L_{H,\Gr},
                 \qquad L_{H,\Gr}-L_H|_{D'}\ \text{effective}.
\end{equation}
Here and below restriction notation includes the pullback
to that resolution when one was needed.

Both sides of \eqref{eq:boundary-line-comparison} are nef.
For the right side this follows from the pure highest-line
metric and \Cref{lem:log-metric}; for the left side it follows
by restricting the original nef line.  Tensor the comparison
by $e_Q$ and add the same compact-mark line on $D'$.  If
$P,Q$ are the resulting nef lines, then $Q-P$ is effective,
and for an ample $A$ on this $(n-1)$-dimensional variety,
\[
 (Q^k-P^k)A^{n-1-k}
  =(Q-P)\sum_{a=0}^{k-1}Q^{k-1-a}P^a A^{n-1-k}\geq0.
\]
Thus $\nu(P)\leq\nu(Q)$.  The metric lemma and
\eqref{eq:boundary-graded-rank} now give
\[
                    \nu(\mathcal H_H|_{D'})<d_H.
\]
Adding $\mathcal H_Q|_{D'}$, whose numerical dimension is
at most $s_Q$, proves the desired strict drop upstairs.

\smallskip\noindent
\emph{Finite nonidentity local monodromy.}\label{proof:finite-boundary}
Let $\gamma\ne1$ denote the original monodromy.  On the
unipotent level its logarithm is zero, so the pure period
extends ordinarily across a general point of $D'$.  Its
limiting flags and compact marks are fixed by $\gamma$.
Indeed, the local cover is $t=z^e$; the continuation around
$t=0$ acts on its sheets by $z\mapsto e^{2\pi\sqrt{-1}/e}z$.
Equivariance of the extended period and marks then gives
the fixed-point assertion at $z=0$ in common flat frames.
The comparison in \eqref{eq:boundary-line-comparison}, now
with a pure ordinary limit, shows that the numerical
dimension on $D'$ is bounded by the differential rank of
these limiting marks.  That rank is at most $r$: in the
ordinary extending flat frame, all $(r+1)$-minors of the
mark differential vanish on the interior and hence on
its extension, including after restriction to $D'$.

Suppose that rank were $r=d_H+s_Q$.  Since $D'$ dominates
$S_Q$, the limiting $Q$ data comes from the interior of
$S_Q$.  In particular $dj_Q$ on $D'$ has rank $s_Q$ and
kernel $\ker(dq|_{D'})$.  Equality of the total rank forces
the $H$ flags along a local general fiber of $q|_{D'}$ to
fill an open subset of $J_H$.  The $Q$ flags are fixed
on this fiber.

The rational automorphism $\gamma$ normalizes the connected
group, and acts on the product $J$ factorwise up to
permutation of simple factors.  An open subset with all
$H$ flag coordinates varying independently cannot be fixed
if any visible $H$ coordinate is moved to a different
coordinate.  If the destination were a $Q$ coordinate it
would be constant; if it were another $H$ coordinate the
fixed-point equations would impose a relation between two
independent coordinates.  Thus every visible complex
$H$ factor is preserved and $\gamma$ acts identically on
its whole flag variety.  Equivariance and faithfulness of
the simple adjoint action imply identity on its Lie
algebra.  Rationality then implies identity on every
Galois-conjugate complex factor, hence on all of $H$.
Consequently $\gamma$ preserves $Q$.  Its action on $Q$
fixes a generic lifted full $Q$ period, because $D'$ dominates
$S_Q$.  \Cref{lem:generic-lie-stabilizer} gives identity on
$Q$ also, a contradiction.  The limiting rank is strictly
less than $r$ in this case as well.

Pullback by the generically finite map $D'\to D_i$ preserves
the numerical dimension of a nef line, so the strict drop
descends to $D_i$.  Finally, if $r=0$, a local fiber of $j$
is open in the base, so its Zariski closure is the entire
base and $H=G$.  Orbit filling gives $\dim J=0$.
Visibility on every retained rational factor forces
$G=1$; the full projected adjoint action then has no
nonidentity boundary monodromy.  Thus $D_0$ is empty.
\end{proof}

\subsection{Removing the marked boundary}

\begin{lemma}[Section estimate]\label{lem:boundary-remove}
Let $M$ be a smooth projective $n$-fold, let $H$ be a nef
rational line, and let $T=\sum_i T_i$ be a reduced divisor
with smooth components.  Suppose $r=\nu(H)>0$ and
$\nu(H|_{T_i})<r$ for every $i$.  If $A$ is ample and
$e>0$ is an integer, then
\[
                         A+tH-eT
\]
is big for all sufficiently large divisible integers $t$.
\end{lemma}

\begin{proof}
Choose fixed ample integral lines $A_i$ on $T_i$ sufficiently
positive that all successive restriction terms obtained
while imposing order $le$ along each $T_j$ inject into
\[
 H^0\bigl(T_i,l(A|_{T_i}+tH|_{T_i}+A_i)\bigr).
\]
This choice is uniform in $l$ and $t$.  Explicitly, choose
ample $C_{ij}$ so that $C_{ij}$ and $C_{ij}+T_j|_{T_i}$
have nonzero sections, and put $A_i=e\sum_j C_{ij}$,
enlarging it if needed.  For $0\leq n_j\leq el$ the line
\[
 lA_i+\sum_j n_jT_j|_{T_i}
 =\sum_j\bigl((el-n_j)C_{ij}
                  +n_j(C_{ij}+T_j|_{T_i})\bigr)
\]
has a nonzero section.  Multiplication by it supplies
exactly the asserted injections, including the normal
twists when $j=i$.

The restriction exact sequences, applied one copy of a
component at a time, give
\begin{align*}
 h^0\bigl(M,l(A+tH-eT)\bigr)
 &\geq h^0\bigl(M,l(A+tH)\bigr)\\
 &\quad-le\sum_i
 h^0\bigl(T_i,l(A|_{T_i}+tH|_{T_i}+A_i)\bigr).
\end{align*}
For each fixed divisible $t$, all the positive lines in
this formula are ample.  Their Hilbert polynomials give,
in the coefficient of $l^n$, a positive supply
\[
                         \frac{(A+tH)^n}{n!}
\]
and a loss at most
\[
 \frac{e}{(n-1)!}\sum_i
                (A|_{T_i}+A_i+tH|_{T_i})^{n-1}.
\]
The supply is a polynomial in $t$ of degree exactly $r$
with positive leading coefficient.  Each term in the
loss has degree at most $r-1$.  For $n=1$ the restriction
terms are simply section dimensions on points and the
same statement holds.  Choose $t$ sufficiently large
that the difference of these leading coefficients is
positive, and then let $l\to\infty$.  This proves bigness.
No asymptotic estimate uniform in both $l$ and $t$ is used.
\end{proof}

\begin{proposition}[Adjoint bigness]\label{prop:adjoint-bigness}
If $\dim S>0$, then $K_S+b\mathcal H$ is big for all
sufficiently large rational $b$.
\end{proposition}

\begin{proof}
By \Cref{lem:initial-adjoint},
$K_S+D_0+c\mathcal H_{\rm u}$ is big for sufficiently
large $c$.  Since $\mathcal H-\mathcal H_{\rm u}$ is nef,
$K_S+D_0+c\mathcal H$ is big.  Choose an integer $e>0$
clearing denominators and an ample integral $A$ with
\[
                   e(K_S+c\mathcal H+D_0)-A
\]
effective.  \Cref{prop:boundary-drop,lem:boundary-remove}
show that $A+t\mathcal H-eD_0$ is big for large divisible
$t$.  Adding the displayed effective divisor proves
bigness of $eK_S+(ec+t)\mathcal H$.  When $\nu(\mathcal H)=0$,
\Cref{prop:boundary-drop} says $D_0=0$, and the initial
bigness already gives the conclusion.  Finally the sum
of a big line and a nef line is big, so every larger
rational $b$ works as well.
\end{proof}

We have proved adjoint bigness on the parameter space.  To use it for a
chosen vector, we need a nef line generated by its Higgs images whose
numerical dimension does not increase on adding the marked line.
The next proposition supplies this second input to the removal of the
period twist in \Cref{sec:nonvanishing}.

\subsection{The rank of the Higgs-tail line}

\begin{proposition}[Higgs tails]\label{prop:tail-rank}
Let $h:W\to S$ be any dominant morphism of smooth
connected projective varieties with connected generic
fiber.  On a dense algebraic open of $W$, suppose that
$\ell\subset h^*E_0$ is an algebraic line subbundle.
Require that the pointwise instability-parabolic
construction of \Cref{lem:marked-flags}, applied to
$\ell$, gives exactly the pulled-back compact marks
$h^*j_{\rm u}$.  This compatibility is part of the
hypothesis.

Take a smooth connected unipotent level
$\pi:\widehat W\to W$ dominating the level used to
define $\mathcal H$, and resolve its boundary and the
generated subspace maps.  Over the variation locus,
let $G_i$ be the span of all contractions of
$\theta^i(\ell)$ in
$\pi^*h^*\Gr_F^{p_*-i}\mathbb A$, with $G_0=\ell$.
Use the resulting subbundles of the extended Hodge
grades on $\widehat W$.  Put
\begin{equation}\label{eq:tail-line-definition}
 U_j=\bigoplus_{i\geq j}G_i,
 \qquad
 P=-\sum_{j\geq0}\det U_j
   =-\sum_{i\geq0}(i+1)\det G_i.
\end{equation}
Only finitely many summands are nonzero.  Then $P$ is
nef and
\begin{equation}\label{eq:tail-rank-equality}
              \nu\bigl(P+\pi^*h^*\mathcal H\bigr)=\nu(P).
\end{equation}
The assertion applies to every such $h$ and every such
compatible line; $\ell$ need not be the original vector
used to construct $S$.
\end{proposition}

The curvature argument will show that a null direction for $P$ fixes
the chosen line and its compact marks.  Passing from that line to the
whole highest Hodge space requires the following representation lemma.
We state it here, use it in the proof of the proposition,
and give its root-theoretic proof afterward.

\begin{lemma}[Degree-one roots]\label{lem:degree-one-root}
Let $\mathfrak g=\bigoplus_k\mathfrak g_k$ be a complex
semisimple Lie algebra with an integral grading and an
adjoint operation $*$ carrying $\mathfrak g_k$ to
$\mathfrak g_{-k}$, arising from a compatible compact
form.  Let $\mathfrak h$ be a graded ideal stable under
$*$, and let $V$ be an irreducible finite-dimensional
representation with compatible grading and highest
graded piece $E$.  The metric on $V$ is chosen so that
representation adjoints are the actions of Lie adjoints.
For $0\ne s\in E$ and $x\in\mathfrak h_1$, suppose
\begin{equation}\label{eq:root-annihilation-hypothesis}
                         [x,\mathfrak h_{-1}]s=0.
\end{equation}
Then $x^*E=0$.
\end{lemma}

\begin{proof}[Proof of Proposition~\ref{prop:tail-rank}]
The connected-fiber hypothesis ensures that, after
shrinking to a smooth locally trivial locus, the
monodromy image of the pullback is the same as that
on $S$.  Passing to its connected level therefore
leaves Zariski-dense monodromy $G$.  The variation
and the orbit calculation of \Cref{lem:orbit-rank}
apply on $\widehat W$.

Each $U_j$ is Higgs invariant because
$\theta G_i\subseteq G_{i+1}\otimes\Omega^1$; this
follows from Higgs integrability and the definition
by all contractions.  On the open locus the Hodge
decomposition is orthogonal.  Write $\mathcal D$ for
its Chern connection.  For a tangent vector $v$ put
$A=\theta(v)$, let $\Pi_j$ be orthogonal projection
onto $U_j$, and set
\[
             B_j(v)=(1-\Pi_j)\mathcal D'_v|_{U_j}.
\]
The Hodge curvature equation and the subbundle
curvature formula give, in a common positive
normalization,
\begin{equation}\label{eq:tail-curvature}
 \Theta_{-\det U_j}(v,\bar v)
   =\bigl\|(1-\Pi_j)A^*|_{U_j}\bigr\|_{\rm HS}^2
                      +\|B_j(v)\|_{\rm HS}^2.
\end{equation}
To see the sign directly, use $A U_j\subseteq U_j$
and write $A$ in upper block-triangular form for
$U_j\oplus U_j^\perp$.  The traces of the two
internal products cancel.  The remaining trace is
minus the square of the off-diagonal block of
$A^*|_{U_j}$ for $\det U_j$, and subbundle curvature
subtracts the square of $B_j(v)$.  Negating the
determinant gives \eqref{eq:tail-curvature}.

The metrics on the extended subbundles have the
two-sided logarithmic norm estimates needed in
\Cref{lem:log-metric}.  Indeed they are restrictions
of the extended Hodge metrics; after resolution
their frames have full rank in the extended Hodge
grades, so the dual estimates also restrict.
Thus every $-\det U_j$, and hence $P$, is nef.

Suppose now that $\Theta_P(v,\bar v)=0$.  Every
nonnegative summand in \eqref{eq:tail-curvature}
vanishes.  Applying its first term for $U_i$ to
$G_i$, the vector $A^*G_i$ belongs to the preceding
Hodge grade, which is orthogonal to $U_i$.  Hence
\begin{equation}\label{eq:tail-zero-actions}
                A^*G_i=0\quad\hbox{for every }i.
\end{equation}
Since $AG_i\subseteq G_{i+1}$, for $g\in G_i$ one
then has $\|Ag\|^2=\langle g,A^*Ag\rangle=0$.
Thus $AG_i=0$ as well.  Moreover $B_0(v)=0$ and
$\mathcal D'$ preserves the Hodge grading, so
$\mathcal D'_v$ preserves $G_0=\ell$.  The flat
connection satisfies
$\nabla^{1,0}=\mathcal D'+\theta$ in the smooth
Hodge splitting.  Its projective derivative on
$\ell$ is consequently zero.  In other words,
\begin{equation}\label{eq:tail-null-line}
                         v\in\ker d\ell.
\end{equation}
The compact marks, which by hypothesis are the
same pointwise functions of $\ell$ as before,
are constant on a local $\ell$-fiber.  At its
constant-rank points \eqref{eq:tail-null-line}
also gives $dj_{\rm u}(v)=0$.

It remains to prove $dE_0(v)=0$.  Apply the line
version of \Cref{lem:orbit-rank} to a local
$\ell$-fiber $T$, with Zariski closure $Z$ and
decomposition $G=H\times Q$.  Work at one of the
smooth generic points $z$ supplied by that lemma.
The preceding calculation places every
$v\in\ker\Theta_{P,z}$ in $T_zT$; fix this same
$T$, $Z$, and $H$ for all such vectors.  The full
$Q$ period is constant on $Z$.  Hence the
infinitesimal period element of every such $v$ satisfies
\[
                    t(v)\in\mathfrak h^{-1,1}.
\]
Write $\ell_z=\C s$.  Evaluating the tangent map
of $\ell$ at $s$ and projecting to the next Hodge
grade sends $d\ell(T_zZ)$ into $G_1$: its flat
derivative is $\mathcal D'_w s+\theta(w)s$, and
only the second term contributes to that grade.
Orbit filling in \eqref{eq:boundary-orbit-fill}
therefore gives
\begin{equation}\label{eq:tail-orbit-action}
                         \mathfrak h^{-1,1}s\subseteq G_1.
\end{equation}
Use the metric supplied by the adjoint tensor
construction, so that representation adjoints
agree with Lie adjoints.  Set $x=t(v)^*$.
By \eqref{eq:tail-zero-actions},
$xG_1=0$.  Also $xs=0$ because $s$ is in the
highest Hodge grade and $x$ raises that grade.
It follows from \eqref{eq:tail-orbit-action} that
\[
                         [x,\mathfrak h^{-1,1}]s=0.
\]
\Cref{lem:degree-one-root}, with
$\mathfrak g_k=\mathfrak g^{k,-k}$, gives
$t(v)E_0=0$.  The flat-chart differential of the
whole highest step is precisely Higgs evaluation
on $E_0$, so $dE_0(v)=0$.

We have proved, at the points where the orbit
calculation applies, the kernel inclusion
\begin{equation}\label{eq:tail-kernel-inclusion}
                  \ker\Theta_P\subseteq
                          \ker\Theta_{\pi^*h^*\mathcal H}.
\end{equation}
For completeness, these points can be chosen to
compute the global curvature ranks.  The Hodge
metrics and subspace data are real analytic on
the variation locus.  Choose the initial point
in the open maximum-rank loci of $\Theta_P$ and
$\Theta_P+\Theta_{\pi^*h^*\mathcal H}$, as well
as the finitely many relevant period, flag, and
projective-line rank loci and the Hodge-generic locus.  Moving
slightly within its $\ell$-fiber reaches the
smooth Zariski-open locus of $Z$ required by
\Cref{lem:orbit-rank}, while staying in the two
curvature-rank opens.  Indeed an open subgerm of
that fiber is Zariski dense in $Z$.  The argument
then establishes \eqref{eq:tail-kernel-inclusion}
at a point where both curvature ranks are
maximal.  Semipositivity implies equality of
those ranks.  Apply \Cref{lem:log-metric} to
$P$ and $P+\pi^*h^*\mathcal H$ to obtain
\eqref{eq:tail-rank-equality}.
\end{proof}

\begin{proof}[Proof of Lemma~\ref{lem:degree-one-root}]
Choose a Cartan in the compact form adapted to the grading,
and choose positive roots so that each simple root has
nonnegative degree.  Metric adjoints exchange opposite
root spaces.  Let $\Delta_0$ be the degree-zero simple
roots, and let $v_\lambda$ be a highest-weight vector
for $V$.  Write $\mathfrak g_0^-$ for the sum of the
negative root spaces of degree zero.  The highest graded piece is
\begin{equation}\label{eq:root-top-module}
                  E=U(\mathfrak g_0^-)v_\lambda.
\end{equation}
In particular it is irreducible for the full degree-zero
reductive algebra.  Indeed, reaching the extremal grading
from the highest weight allows only degree-zero
lowerings; the highest-weight theory for that Levi
algebra then gives \eqref{eq:root-top-module}.

Every degree-one positive root contains exactly one
degree-one simple root $\alpha$, with coefficient one,
and otherwise only simple roots of degree zero.  Let
$\mathfrak b^+_\alpha$ be the sum of its root spaces and
$\mathfrak b^-_\alpha$ the sum of the opposite root
spaces.  These are $\mathfrak g_0$-modules, and $*$
exchanges them.  Roots of degree one in $\mathfrak h$
are exactly the blocks belonging to its simple factors.
For $\alpha\ne\beta$,
\begin{equation}\label{eq:root-cross-bracket}
                 [\mathfrak b^+_\alpha,
                         \mathfrak b^-_\beta]=0.
\end{equation}
Indeed, the difference of the relevant roots has a
positive coefficient at $\alpha$ and a negative one at
$\beta$, and hence is neither a root nor zero.

Write $x=\sum_\alpha x_\alpha$ using these blocks in
$\mathfrak h_1$.  For every nonzero $x_\alpha$,
\eqref{eq:root-annihilation-hypothesis} and
\eqref{eq:root-cross-bracket} imply
$[x_\alpha,\mathfrak b^-_\alpha]s=0$.  Hence the incidence
\[
 I_\alpha=
 \left\{([z],[u])\in
  \mathbf P(\mathfrak b^+_\alpha)\times\mathbf P(E):
             [z,y]u=0\text{ for every }y\in\mathfrak b^-_\alpha\right\}
\]
is nonempty, closed, and projective.  It is stable under
the full connected degree-zero group.  Apply the Borel
fixed-point theorem to the Borel of this group determined
by the chosen positive roots and containing the full Cartan.
The theorem used here says that a
connected solvable algebraic group acting on a nonempty
complete variety has a fixed point
\cite[Corollary~17.3]{Milne}.  A fixed pair in $I_\alpha$
has the form $([e_\rho],[v_\lambda])$: the first
coordinate is a root line since the full Cartan acts,
and the second is the unique highest line of the
irreducible degree-zero module $E$.  Irreducibility of
$\mathfrak b^+_\alpha$ is not required.

Let $C_\alpha$ be the connected component containing
$\alpha$ in the Dynkin subdiagram on
$\Delta_0\cup\{\alpha\}$.  Connectedness of root support
gives $\operatorname{Supp}(\rho)\subseteq C_\alpha$.
If the inclusion were strict, there would be a missing
zero-degree node $\delta$ adjacent to the support.
Its coefficient in $\rho$ is zero, while some adjacent
support coefficient is positive, so
$\langle\rho,\delta^\vee\rangle<0$.  The root-string
property implies that $\rho+\delta$ is a root.  Thus
$e_\delta$ does not preserve the line $\C e_\rho$,
contrary to the Borel-fixed condition.  Therefore
\[
                     \operatorname{Supp}(\rho)=C_\alpha.
\]
Taking $y=e_{-\rho}$ in the incidence equation gives
$\langle\lambda,\rho^\vee\rangle=0$.  The simple-coroot
expansion of $\rho^\vee$ has strictly positive
coefficients on its support.  Dominance of $\lambda$
therefore gives
$\langle\lambda,\delta^\vee\rangle=0$ for every
$\delta\in C_\alpha$.  Every root indexing
$\mathfrak b^+_\alpha$ has support in $C_\alpha$, so
its coroot also pairs trivially with $\lambda$.
Finite-dimensional root $\mathfrak{sl}_2$ theory implies
that every vector in $\mathfrak b^-_\alpha$ kills
$v_\lambda$.  Since $\mathfrak b^-_\alpha$ is a
$\mathfrak g_0$-module, commuting it past the
degree-zero lowerings in \eqref{eq:root-top-module}
shows that it kills all of $E$.  These root-string and
highest-weight facts are the usual semisimple
representation theory; see \cite[\S\S9--10,20--21]{Humphreys}.
Finally $x_\alpha^*\in\mathfrak b^-_\alpha$ for every
nonzero component, and summing proves $x^*E=0$.
\end{proof}
